%% kapa: exam=1 round=26 year=2015 period=1 form=A subject=economics subject_ko=경제원론 seq=17 number=57
%% source: Q-Net 감정평가사 기출문제 (cst00309) artlSeq=5138001 file=extracted/1008/5138001/제26회_감정평가사 1차_기출/제26회_감정평가사 1차_기출/1교시_감평_A형.pdf
%% conversion: kordoc 4.24.51 -> tools/gen1.py (auto); review=proofread-sonnet; flags=none
\begin{question}{57}
다음 표는 소비의 배제성과 경합성의 존재 유무에 따라 재화를 분류하고 있다. 다음 표에서 $C$에 해당하는 재화로 옳은 것은?
\par{\small\centering\begin{tabular}{cc|c|c|}
\multicolumn{2}{c}{} & \multicolumn{2}{c}{경합성} \\
\multicolumn{2}{c}{} & 있음 & 없음 \\ \cline{3-4}
\multirow{2}{*}{배제성} & 있음 & $A$ & $B$ \\ \cline{3-4}
 & 없음 & $C$ & $D$ \\ \cline{3-4}
\end{tabular}\par}
\begin{choices}
  \item 사적(私的) 재화
  \item 유료 도로
  \item 국방서비스
  \item 유료 케이블 TV
  \item 공해(公海) 상의 물고기
\end{choices}
\end{question}
\answer{5}
